\documentclass[aspectratio=169,11pt]{beamer}

\usetheme{metropolis}
\usepackage{graphicx}
\usepackage{booktabs}
\usepackage{amsmath}
\usepackage{listings}
\usepackage{xcolor}
\usepackage{hyperref}

% Colors
\definecolor{codegreen}{rgb}{0,0.6,0}
\definecolor{codegray}{rgb}{0.5,0.5,0.5}
\definecolor{codebg}{rgb}{0.97,0.97,0.97}

% Code listing style
\lstset{
  backgroundcolor=\color{codebg},
  basicstyle=\ttfamily\scriptsize,
  breaklines=true,
  captionpos=b,
  commentstyle=\color{codegreen},
  keywordstyle=\color{blue},
  numberstyle=\tiny\color{codegray},
  showstringspaces=false,
  frame=single,
  framerule=0pt,
}

% Metadata — fill these in per PI check-in
\newcommand{\cyclerange}{XX--YY}       % e.g. 01--04
\newcommand{\currentcycle}{YY}          % current cycle number
\newcommand{\projecttitle}{Project Title}
\newcommand{\directionval}{--}
\newcommand{\velocityval}{--}

\title{PI Check-in — Cycles \cyclerange}
\subtitle{\projecttitle}
\date{\today}
\author{Claude (PhD Student)}
\institute{Supervised by: Daniel Murnane}

\begin{document}

\maketitle

% ============================================================
% SLIDE 1: Status & Context
% Multi-cycle overview: where we started, where we are now
% ============================================================
\begin{frame}{Status \& Context}
  \begin{columns}[T]
    \begin{column}{0.5\textwidth}
      \textbf{Cycles covered:} \cyclerange \\
      \textbf{Current cycle:} \currentcycle \\
      \textbf{Direction:} \directionval\% \\
      \textbf{Velocity:} \velocityval\%
    \end{column}
    \begin{column}{0.5\textwidth}
      \textbf{Direction/Velocity trajectory:} \\
      % Show how d/v evolved across cycles
      % e.g. a small table or sparkline
      \begin{itemize}
        \item Cycle XX: d=??, v=??
        \item Cycle YY: d=??, v=??
      \end{itemize}
    \end{column}
  \end{columns}
\end{frame}

% ============================================================
% SLIDE 2: Summary of Work
% What was accomplished across N cycles — key milestones
% ============================================================
\begin{frame}{Summary of Work (Cycles \cyclerange)}
  \begin{itemize}
    \item TODO: Key milestone from cycle XX
    \item TODO: Key milestone from cycle XX+1
    \item TODO: Key milestone from cycle YY
  \end{itemize}
\end{frame}

% ============================================================
% SLIDE 3: Cumulative Results
% Trajectory of metrics, best results, comparisons across cycles
% ============================================================
\begin{frame}{Cumulative Results}
  % \includegraphics[width=\textwidth]{../results/metric_trajectory.pdf}
  TODO: plots showing metric evolution across cycles, comparison tables,
  best results achieved so far
\end{frame}

% ============================================================
% SLIDE 4: Postdoc's Assessment
% This content is drawn from .postdoc/pi_brief.md
% The enforce_checkin hook injects it as meeting context
% Summarize the postdoc's key observations here
% ============================================================
\begin{frame}{Postdoc's Assessment}
  % Populated from the postdoc's PI brief
  \begin{itemize}
    \item TODO: Postdoc's overall assessment
    \item TODO: Key concerns raised
    \item TODO: Patterns observed
  \end{itemize}
\end{frame}

% ============================================================
% SLIDE 5: Questions for PI
% Accumulated from both student and postdoc across cycles
% ============================================================
\begin{frame}{Questions for PI}
  \begin{enumerate}
    \item TODO: Strategic question accumulated over cycles
    \item TODO: Decision that needs PI input
    \item TODO: Postdoc-flagged question
  \end{enumerate}
\end{frame}

% ============================================================
% SLIDE 6: Proposed Direction
% Next period's plan and direction/velocity proposal
% ============================================================
\begin{frame}{Proposed Direction}
  \textbf{Plan for next period (cycles YY--ZZ):}
  \begin{itemize}
    \item TODO: High-level research direction
  \end{itemize}

  \vspace{1em}

  \textbf{Proposed direction/velocity:} \\
  Direction: \directionval\% $\rightarrow$ ??\% \\
  Velocity: \velocityval\% $\rightarrow$ ??\%
\end{frame}

\end{document}
